\section{Cinemática diferencial do efetuador e composições de orientação}


Para representar a atitude (orientação) de um corpo rígido, neste caso, o manipulador robótico, é necessário detalhar a relação entre os \emph{frames (ou sistema de coordenadas)} e um sistema de referência global.

\subsection{Relação entre frames e sistema de referência global}
Como existem múltiplos
\emph{frames} envolvidos, como o Inercial, centrado na Terra $\{I\}$;
o \emph{frame} local da espaçonave, $\{LVLH\}$, que por sua vez está no centro de massa do \emph{chaser};
e o \emph{corpo $\{B\}$}, fixado no centro de massa do manipulador robótico,
se faz necessário definir a relação entre esses \emph{frames} e o sistema de referência global.
A matriz de rotação é utilizada para descrever as coordenadas de um vetor de um sistema para outro \cite{craig2005}.

\begin{equation}
\label{eq:LocalparaInercial}
{}^{I}\mathbf{R}_{B} = {}^{I}\mathbf{R}_{LVLH}\;{}^{LVLH}\mathbf{R}_{B}
\end{equation}

Enquanto que a matriz de rotação pode ser descrita por:
\begin{equation}
\label{eq:LocalparaInercial_yawpitchroll}
{}^{LVLH}\mathbf{R}_{B} (\psi,\theta,\phi) = \mathbf{R}_z(\psi)\,\mathbf{R}_y(\theta)\,\mathbf{R}_x(\phi)
\end{equation}

A composição de \eqref{eq:LocalparaInercial_yawpitchroll} corresponde à sequência extrínseca ZYX adotada anteriormente, e para melhor compreensão dos termos nesta seção, os símbolos de rotação em cada eixo principal foram definidos na tabela \eqref{tab:euler_LVLH}.

\begin{table}[ht]
    \centering
    \footnotesize
    \caption{Relação entre ângulos de Euler e eixos do sistema LVLH.}
    \label{tab:euler_LVLH}
    \begin{tabular}{lll}
        \hline
        \textbf{Ângulo} & \textbf{Eixo de rotação} & \textbf{Descrição} \\
        \hline
        $\phi$ (roll, rolagem)
        & $x$ (radial)
        & Rotação em torno do eixo radial \\

        $\theta$ (pitch, arremesso)
        & $y$ (along-track)
        & Rotação em torno do eixo tangente à órbita \\

        $\psi$ (yaw, guinada)
        & $z$ (cross-track)
        & Rotação lateral em torno do eixo normal ao plano orbital \\

        \hline
    \end{tabular}
\end{table}

Como o manipulador robótico deste documento possui 5 juntas, pode-se adaptar \eqref{eq:produtorio_generico} para essa totalidade de juntas, como indicado a seguir:
\begin{equation}
{}^{0}\!T_E(q)=
A_1(\theta_1)
A_2(\theta_2)
A_3(\theta_3)
A_4(\theta_4)
A_5(d_5),
\end{equation}
onde $A_i$ são as transformações elementares da tabela DH, e $\theta_i$ são os ângulos revolutivos das juntas, e $d_5$ é o deslocamento da junta prismática. A matriz homogênea ${}^0T_E$ que descreve tanto a orientação como a posição do efetuador em relação à base é:
\begin{equation}
{}^0T_E(q) \;=\;
\begin{bmatrix}
r_{11} & r_{12} & r_{13} & p_x \\
r_{21} & r_{22} & r_{23} & p_y \\
r_{31} & r_{32} & r_{33} & p_z \\
0      & 0      & 0      & 1
\end{bmatrix}.
\label{eq:T0E_completa}
\end{equation}

A submatriz de rotação $\mathbf{R}_E(q)$ e o vetor de posição do efetuador $\vec{p}_E(q)$ são, respectivamente:
\begin{equation}
\mathbf{R}_E(q) \;=\;
\begin{bmatrix}
r_{11} & r_{12} & r_{13} \\
r_{21} & r_{22} & r_{23} \\
r_{31} & r_{32} & r_{33}
\end{bmatrix},
\label{eq:RE_def}
\end{equation}
\begin{equation}
\vec{p}_E(q) \;=\;
\begin{bmatrix}
p_x \\ p_y \\ p_z
\end{bmatrix}.
\label{eq:pE_def}
\end{equation}






